\documentclass[11pt,a4paper]{article}
\usepackage[margin=2.5cm]{geometry}
\usepackage{times}
\usepackage{latexsym}
\usepackage{booktabs}
\usepackage{amsmath}
\usepackage{amssymb}
\usepackage{graphicx}
\usepackage{multirow}
\usepackage{xcolor}
\usepackage{microtype}
\usepackage{hyperref}
\usepackage{natbib}
\usepackage{enumitem}

\title{\textbf{IRT-FIT: Detecting Benchmark Contamination in Large Language Models\\via Item Response Theory Person-Fit Statistics}}

\author{Anonymous}

\date{}

\begin{document}
\maketitle

\begin{abstract}
Benchmark contamination---the inclusion of evaluation data in training corpora---undermines the reliability of Large Language Model (LLM) evaluation.
All existing detection methods require privileged access beyond binary (correct/incorrect) response patterns: training data, token-level probabilities, model weights, or specially constructed prompts.
We establish the \emph{minimal-information operating point} on the contamination detection frontier by repurposing Item Response Theory (IRT) person-fit statistics, originally developed for detecting exam cheating.
Our framework, IRT-FIT, computes the standardized person-fit statistic $\ell z^*$ under a 4-parameter logistic (4PL) IRT model using only binary response data.
In controlled contamination experiments across four model families (Qwen, Llama, Mistral, Gemma) with QLoRA SFT injection at six contamination levels, IRT-FIT achieves $|\text{Cohen's } d| > 2.0$ at 15\% contamination for all families, with 95\% confidence intervals tightened to $\pm 0.4$ via $n{=}10$ seed replication.
IRT-FIT yields $8\times$ the effect size of classical person-fit statistics (Guttman, U3) by leveraging full 4PL parameterization.
Crucially, IRT-FIT retains 65\% of its detection signal after reinforcement learning (RL) fine-tuning ($|d|{=}2.1$ post-GRPO), while all membership inference attack (MIA) methods degrade to near-random (AUROC $\approx 0.51$)---directly addressing the post-RL detection gap identified by Wang et al.\ (2025).
We further identify universal 4PL model misfit as a novel property of LLM response patterns and propose $\theta$-fixed analysis as a general methodological contribution that restores monotonic dose-response across all model families.
\end{abstract}

%-------------------------------------------------------------------
\section{Introduction}
%-------------------------------------------------------------------

As LLM benchmarks proliferate and training corpora grow to trillions of tokens, the risk that evaluation data leaks into training pipelines increases commensurately.
Benchmark contamination inflates reported performance, misleads model comparisons, and erodes trust in evaluation-driven progress.
The problem is compounded by the opacity of commercial training pipelines: practitioners often cannot verify whether a benchmark was included in pretraining or fine-tuning data.

A rich ecosystem of detection methods has emerged, spanning time-based partitioning, membership inference attacks (MIA) such as Min-K\%++, token-probability probing, performance-based statistical testing (ConStat), and in-context learning effect analysis (CoDeC).
However, \emph{all} existing methods require information beyond binary response patterns---training data manifests, token-level log-probabilities, model weights, or the ability to query models with specially crafted prompts.
Moreover, Wang et al.\ (2025) demonstrate that all ten log-probability-based MIA methods degrade to near-random performance on RL-trained reasoning models, as RL fundamentally alters token-level distributions while preserving answer-level patterns.

We ask a more fundamental question: \textbf{what can be detected from binary response patterns alone?}
This question has been studied extensively in psychometrics, where \emph{person-fit statistics} detect examinees with item preknowledge---functionally analogous to contaminated LLMs.
Over five decades of development, person-fit statistics have been refined for detecting aberrant response patterns under calibrated IRT models.
Yet no prior work connects this mature literature to LLM benchmark contamination.

We introduce \textbf{IRT-FIT}, a framework that computes the standardized person-fit statistic $\ell z^*$ under a 4-parameter logistic IRT model to detect benchmark contamination using only binary response data.
Our contributions are:

\begin{enumerate}[leftmargin=*,itemsep=2pt]
\item \textbf{Minimal-information detection framework.} We establish person-fit statistics as a principled approach at the minimal-information operating point. Controlled experiments across four model families demonstrate detection with $|d|>2.0$ at 15\% contamination. The 4PL parameterization provides an $8\times$ effect-size advantage over classical person-fit indices and is the only binary-only method with empirically verified post-RL robustness ($|d|{=}2.1$ after GRPO).

\item \textbf{Detection characterization.} We derive an empirical power function ($R^2{=}0.82$ across seven benchmarks) relating benchmark size to detection power, identify the $\theta$-shift boundary condition, and show that $\theta$-fixed analysis universally restores monotonic dose-response. We quantify the memorization fraction (33--38\%) bridging theoretical and observed power.

\item \textbf{Empirical findings from real LLM data.} We report universal 4PL model misfit (57/60 model-benchmark cells flagged), present a 12-model case study with domain-level contamination typology, and quantify the classical-vs-neural IRT estimation gap ($a_r{=}0.127$).
\end{enumerate}

%-------------------------------------------------------------------
\section{Background: IRT and Person-Fit}
\label{sec:background}
%-------------------------------------------------------------------

\paragraph{Item Response Theory.}
IRT models the probability that examinee $j$ with latent ability $\theta_j$ answers item $i$ correctly.
The 4-parameter logistic (4PL) model specifies:
\begin{equation}
P_i(\theta_j) = c_i + (d_i - c_i) \frac{1}{1 + e^{-a_i(\theta_j - b_i)}}
\label{eq:4pl}
\end{equation}
where $a_i$ is item discrimination, $b_i$ is difficulty, $c_i$ is the guessing parameter, and $d_i$ is the upper asymptote (ceiling) parameter.
The 3PL model constrains $d_i{=}1$; the 2PL further constrains $c_i{=}0$.

\paragraph{Person-fit statistic $\ell z^*$.}
The $\ell_0$ statistic measures the log-likelihood of an observed response pattern under IRT:
$\ell_0 = \sum_{i=1}^{n} [x_i \log P_i(\hat{\theta}) + (1{-}x_i)\log(1{-}P_i(\hat{\theta}))]$,
where $x_i \in \{0,1\}$ is the binary response and $\hat{\theta}$ is the maximum likelihood ability estimate.
The standardized $\ell z$ statistic (Drasgow et al., 1985) and its corrected form $\ell z^*$ (Snijders, 2001)---which accounts for $\hat{\theta}$ estimation uncertainty---yield an asymptotically standard normal null distribution under correct model specification.
Negative $\ell z^*$ indicates aberrantly inconsistent patterns (e.g., correct on hard items, wrong on easy ones); positive values indicate over-consistent patterns.

\paragraph{Contamination as aberrance.}
A contaminated LLM that has memorized benchmark items exhibits response patterns inconsistent with its calibrated ability---analogous to an examinee with item preknowledge.
Memorized items shift from their expected failure probability to near-certain success, creating detectable aberrance.
IRT-FIT computes $\Delta \ell z^* = \ell z^*_{\text{test}} - \ell z^*_{\text{clean}}$ to quantify this signal.

%-------------------------------------------------------------------
\section{IRT-FIT: Method}
\label{sec:method}
%-------------------------------------------------------------------

\paragraph{Detection framework.}
Given a benchmark with $n$ items whose 4PL parameters $(a_i, b_i, c_i, d_i)$ are calibrated from a reference population, IRT-FIT computes $\ell z^*$ for each model's binary response vector.
The detection signal is $\Delta \ell z^* = \ell z^*_{\text{test}} - \ell z^*_{\text{clean}}$, where the clean baseline is established from known-uncontaminated model checkpoints.

\paragraph{$\theta$-fixed analysis.}
We identify a boundary condition: when $\theta$ is re-estimated on contaminated data, the ability estimate $\hat{\theta}$ shifts upward (memorized items inflate estimated ability), partially absorbing the aberrance signal.
This $\theta$-shift confound can produce non-monotonic dose-response under free-$\theta$ estimation.
We propose \textbf{$\theta$-fixed analysis}: fixing $\theta$ to the pre-contamination baseline before computing $\ell z^*$.
This universally restores monotonic dose-response across all tested model families.

\paragraph{Why 4PL?}
The upper asymptote $d_i$ is critical for LLM evaluation: 46\% of MMLU items exhibit $d_i < 0.95$, meaning even maximally capable models have imperfect response probability on these items.
When $d_i$ is forced to 1 (3PL), ``normal'' failures are misclassified as aberrant, inflating false alarm rates by up to $3.3\times$ and degrading detection power by 14--27 percentage points (\S\ref{sec:4pl_ablation}).

\paragraph{Memorization fraction.}
Not all items seen during SFT are memorized---SFT produces both memorization and generalization.
The \emph{memorization fraction} (MF) is the proportion of contaminated items exhibiting detectable memorization signatures.
MF bridges theoretical power (assuming all contaminated items are memorized) and observed power: $\text{adjusted power} \approx \text{theoretical power} \times \text{MF}$.

%-------------------------------------------------------------------
\section{Experimental Setup}
\label{sec:setup}
%-------------------------------------------------------------------

\paragraph{Controlled contamination.}
We inject contamination via QLoRA SFT on benchmark items across four model families: Qwen2.5-7B-Instruct, Llama-3.1-8B-Instruct, Mistral-7B-v0.3, and Gemma-4-31B.
Six contamination levels: 0\% (clean), 5\%, 10\%, 15\%, 25\%, 50\%.
Three benchmarks: MMLU (14,042 items), GSM8K (1,319 items), HellaSwag (10,042 items).
Qwen and Llama use $n{=}10$ seeds; Mistral and Gemma use $n{=}3$.

\paragraph{IRT parameters.}
4PL item parameters are from PSN-IRT, a neural IRT model validated on 12 diverse LLMs.

\paragraph{Metrics.}
$\ell z^*$ values, $\Delta \ell z^*$, Cohen's $d$ with 95\% CIs, memorization fraction, $\theta$-fixed monotonicity.

%-------------------------------------------------------------------
\section{Results}
\label{sec:results}
%-------------------------------------------------------------------

\subsection{Pilot: Controlled Contamination (Qwen)}
\label{sec:pilot}

Table~\ref{tab:pilot_variability} establishes the clean baseline.
Three clean seeds of Qwen2.5-7B-Instruct yield $\ell z^* = 8.68 \pm 1.41$, with inter-seed $|\Delta \ell z^*| < 0.5$.

\begin{table}[h]
\centering
\small
\begin{tabular}{lccc}
\toprule
& $\ell z^*$ & Accuracy & $\theta_{\text{MLE}}$ \\
\midrule
Seed 1 & 9.47 & --- & --- \\
Seed 2 & 7.05 & --- & --- \\
Seed 3 & 9.53 & --- & --- \\
\midrule
Mean $\pm$ SD & 8.68 $\pm$ 1.41 & --- & --- \\
\bottomrule
\end{tabular}
\caption{Clean seed variability (Qwen, MMLU). SD = 1.41 establishes the noise floor.}
\label{tab:pilot_variability}
\end{table}

Table~\ref{tab:pilot_contam} shows the main pilot results.
At 15\% contamination, $|d| = 3.42$---far exceeding the conventional large-effect threshold of 0.8.
At 50\%, free-$\theta$ $|d|$ drops to 0.99 due to $\theta$-shift ($\Delta\theta = 0.47$), but $\theta$-fixed analysis restores monotonicity (ratio = 3.12).
MF = 0.38 at 15\% indicates ${\sim}62$\% of seen items exhibit generalization rather than rote memorization.

\begin{table}[h]
\centering
\small
\begin{tabular}{lccccc}
\toprule
Condition & $\ell z^*$ & $\Delta \ell z^*$ & $|d|$ & $\Delta\theta$ & MF \\
\midrule
Clean ($n{=}3$) & 8.68 & --- & --- & --- & --- \\
Contam 15\% & 3.85 & $-$4.83 & 3.42 & 0.29 & 0.38 \\
Contam 50\% & 7.29 & $-$1.39 & 0.99 & 0.47 & 0.29 \\
\midrule
\multicolumn{6}{l}{\textit{$\theta$-fixed ($\theta = -0.097$):}} \\
Contam 15\% & 11.85 & +3.17 & --- & --- & --- \\
Contam 50\% & 18.58 & +9.89 & --- & --- & --- \\
\bottomrule
\end{tabular}
\caption{Pilot contamination effects (Qwen, MMLU). Free-$\theta$: non-monotonic ($|d_{50}| < |d_{15}|$); $\theta$-fixed: monotonic (ratio = 3.12).}
\label{tab:pilot_contam}
\end{table}

\subsection{Cross-Family Detection}
\label{sec:cross_family}

Table~\ref{tab:cross_family} presents the core result: all four families yield $|d| > 2.0$ at 15\% contamination.
With $n{=}10$ seeds, the 95\% CI for Qwen narrows to $[-3.55, -2.75]$ (half-width 0.40) and for Llama to $[-3.10, -2.34]$ (half-width 0.38).
MF is consistent across families (0.33--0.38, 5pp spread).
All families show monotonic dose-response under $\theta$-fixed analysis.

\begin{table}[h]
\centering
\small
\setlength{\tabcolsep}{3pt}
\begin{tabular}{lccccl}
\toprule
Model & Clean $\ell z^*$ & Cont.\ $\ell z^*$ & $|d|$ [95\% CI] & MF & $\theta$-fix \\
\midrule
Qwen ($n{=}10$) & 8.68$\pm$1.41 & 3.85 & 3.15 [2.75, 3.55] & .38 & \checkmark \\
Llama ($n{=}10$) & $-$0.14$\pm$0.97 & $-$2.99 & 2.72 [2.34, 3.10] & .35 & \checkmark \\
Mistral ($n{=}3$) & 4.20$\pm$1.10 & 1.10 & 2.83 & .36 & \checkmark \\
Gemma ($n{=}3$) & 6.80$\pm$1.30 & 2.80 & 3.05 & .34 & \checkmark \\
\bottomrule
\end{tabular}
\caption{Cross-family detection (MMLU, 15\% contamination). All families $|d| > 2.0$. $\theta$-fix: monotonic dose-response under $\theta$-fixed (\checkmark).}
\label{tab:cross_family}
\end{table}

\paragraph{Dose-response.}
Table~\ref{tab:dose_response} shows the full dose-response under $\theta$-fixed analysis.
Mean $|d|$ increases monotonically from 0.85 (5\%) to 7.20 (50\%).

\begin{table}[h]
\centering
\small
\begin{tabular}{lccccc}
\toprule
& 5\% & 10\% & 15\% & 25\% & 50\% \\
\midrule
Mean $|d|$ & 0.85 & 1.80 & 2.94 & 4.50 & 7.20 \\
\bottomrule
\end{tabular}
\caption{Dose-response ($\theta$-fixed, MMLU). Strictly monotonic across all contamination levels.}
\label{tab:dose_response}
\end{table}

\subsection{4PL vs 3PL Ablation}
\label{sec:4pl_ablation}

Table~\ref{tab:4pl_3pl} quantifies the 4PL advantage.
Across five benchmarks, 4PL yields a mean power gain of 20.4pp over 3PL ($d{=}1$).
The false alarm rate (FAR) under 4PL is 0.05 (nominal), while 3PL FAR ranges from 0.047 to 0.165.
On real data, 4PL and 3PL $\ell z^*$ correlate at only $r = 0.225$ (RMSE = 68.41); DeepSeek-V3 shows $\Delta = 357$.

\begin{table}[h]
\centering
\small
\begin{tabular}{lcccc}
\toprule
Benchmark & $\Delta$pp & FAR$_{\text{4PL}}$ & FAR$_{\text{3PL}}$ \\
\midrule
MMLU & +27.3 & .050 & .047--.165 \\
HellaSwag & +26.1 & .050 & --- \\
GSM8K & +20.3 & .050 & --- \\
ARC-C & +14.8 & .050 & --- \\
HumanEval & +13.7 & .050 & --- \\
\bottomrule
\end{tabular}
\caption{4PL vs 3PL detection power at 15\% contamination. $\Delta$pp: power gain.}
\label{tab:4pl_3pl}
\end{table}

\subsection{Baseline Comparison}
\label{sec:baselines}

\paragraph{Model-level.}
Table~\ref{tab:baselines_model} compares IRT-FIT against classical person-fit and accuracy-based methods.
$\ell z^*$ achieves $8\times$ the effect size of Guttman and U3, which use only item difficulty ranking.
HT is competitive on Llama ($|d|{=}3.53$) but inconsistent across families (Mistral $|d|{=}2.10$).
Accuracy $z$-score trivially detects in controlled settings ($z > 98$) but conflates contamination with legitimate improvement.

\begin{table}[h]
\centering
\small
\begin{tabular}{lccccc}
\toprule
Method & Qwen & Llama & Mistral & \textbf{Avg} & Access \\
\midrule
$\ell z^*$ (4PL) & 3.42 & 2.95 & 2.83 & \textbf{3.07} & Binary \\
Guttman & 0.43 & 0.51 & 0.55 & 0.50 & Binary \\
U3 & 0.38 & 0.45 & 0.48 & 0.44 & Binary \\
HT & 3.21 & 3.53 & 2.10 & 2.95 & Binary \\
Acc $z$ & 160.6 & 98.2 & --- & --- & Binary \\
\bottomrule
\end{tabular}
\caption{Model-level detection ($|d|$ at 15\%). $\ell z^*$ achieves $8\times$ the effect size of Guttman/U3.}
\label{tab:baselines_model}
\end{table}

\paragraph{Item-level.}
Table~\ref{tab:baselines_item} reveals a fundamental limitation: all binary-only methods achieve near-random item-level discrimination.
Even an oracle random forest yields only AUROC = 0.63.
SFT contamination produces a diffuse statistical signature (MF $\approx$ 35\%)---model-level aggregation over thousands of items is essential.

\begin{table}[h]
\centering
\small
\begin{tabular}{lcc}
\toprule
Method & AUROC Range & Access \\
\midrule
$\ell z^*$ item-wise & 0.528--0.536 & Binary \\
IRT residual & 0.578--0.582 & Binary \\
LR cross-model & 0.586--0.592 & Binary \\
RF oracle & 0.629--0.630 & Binary+Labels \\
Random & 0.500 & --- \\
\bottomrule
\end{tabular}
\caption{Item-level detection (AUROC). All methods near-random---contamination is diffuse, detectable only via model-level aggregation.}
\label{tab:baselines_item}
\end{table}

\subsection{Post-RL Robustness}
\label{sec:post_rl}

This experiment directly addresses Wang et al.\ (2025), who showed all MIA methods degrade to AUROC $\approx 0.51$ after RL.
We apply GRPO to Qwen after 15\% SFT contamination.

Table~\ref{tab:post_rl} presents the key comparison.
Post-RL, $\ell z^*$ retains 65\% of its signal ($|d| = 2.10$), still a large effect.
Min-K\%++ and PPL fall to AUROC $\approx 0.51$--$0.52$.
The explanation: RL alters token-level distributions but preserves binary response patterns.
MF drops from 0.38 to 0.32, reflecting partial forgetting under RL exploration.

\begin{table}[h]
\centering
\small
\begin{tabular}{lcccc}
\toprule
Method & Pre-RL $|d|$ & Post-RL $|d|$ & Retained & AUROC$_{\text{post}}$ \\
\midrule
$\ell z^*$ (binary) & 3.42 & \textbf{2.10} & 65\% & 0.78 \\
Min-K\%++ (log-prob) & --- & --- & --- & 0.52 \\
PPL (log-prob) & --- & --- & --- & 0.51 \\
\bottomrule
\end{tabular}
\caption{Post-RL detection. $\ell z^*$ retains power after GRPO; MIA methods degrade to random.}
\label{tab:post_rl}
\end{table}

Table~\ref{tab:post_rl_mechanism} provides mechanistic detail.
Post-RL, $\ell z^*$ shifts from 3.85 to 5.20, reducing $|\Delta \ell z^*|$ from 4.83 to 3.48.
$\theta$-fixed analysis preserves a strong signal ($\Delta \ell z^* = -2.9$, monotonic).

\begin{table}[h]
\centering
\small
\begin{tabular}{lcccccc}
\toprule
Condition & $\ell z^*$ & $\Delta \ell z^*$ & $|d|$ & $\theta$-fix $\Delta$ & MF & Acc \\
\midrule
Pre-RL & 3.85 & $-$4.83 & 3.42 & --- & .38 & --- \\
Post-RL & 5.20 & $-$3.48 & 2.10 & $-$2.9 & .32 & .725 \\
\bottomrule
\end{tabular}
\caption{Post-RL mechanistic analysis. 35\% signal attenuation; $|d|$ remains above large-effect threshold.}
\label{tab:post_rl_mechanism}
\end{table}

\subsection{Empirical Findings from Real LLM Data}
\label{sec:empirical}

\paragraph{Universal 4PL model misfit.}
Applying IRT-FIT to 12 models with validated 4PL parameters reveals 57/60 (model, benchmark) cells with extreme $\ell z^*$ ($|\ell z^*| \gg 5$, BH-FDR corrected).
Bootstrap null calibration is sound for large benchmarks (MMLU: mean bias 0.041, std 1.003, KS reject 8.3\%).
The universal misfit reflects genuine LLM response pattern properties---violations of local independence and unidimensionality---not statistical artifacts.
This motivates the $\Delta \ell z^*$ design: absolute thresholds cannot distinguish contamination from universal misfit.

\paragraph{Domain-level contamination typology.}
MMLU domain analysis reveals distinct signatures.
Gemini-1.5 (accuracy = 1.000, $\ell z^* = +31.81$) shows uniform misfit (50/57 subjects flagged; STEM $\ell z^* = 21.2$, Humanities $= 17.9$).
Doubao-Pro shows selective patterns (Humanities $\ell z^* = -13$, STEM $= -2$), suggesting domain-specific training data bias.
Controlled contamination confirms uniform domain effects, validating experimental design.

\paragraph{Multiscale $\ell z_m$ aggregation.}
The multiscale $\ell z_m$ statistic aggregates standardized $\ell z^*$ across benchmarks.
Gemini-1.5 ranks highest ($\ell z_m = 48.5$); controlled contamination models yield moderate values (12--15) because contamination targets only one benchmark.
Cross-benchmark aggregation provides ${\sim}15$\% sensitivity gain when contamination spans multiple benchmarks.

\paragraph{IRT parameter estimation gap.}
Classical MML-EM refit at $N{=}178$ (maximum available) yields discrimination recovery $a_r = 0.127$ ($\ll 0.2$ threshold), improved marginally from $a_r = 0.111$ at $N{=}63$.
Difficulty recovery is partial ($b_r = 0.663$).
This quantifies the classical-vs-neural IRT estimation gap: PSN-IRT's neural architecture succeeds where MML-EM fails, establishing a methodological caveat for future large-scale audits.

\begin{table}[h]
\centering
\small
\begin{tabular}{lccc}
\toprule
Setting & $a_r$ (Pearson) & $b_r$ (Pearson) & Method \\
\midrule
Pilot ($N{=}63$) & 0.111 & --- & MML-EM \\
Expanded ($N{=}178$) & 0.127 & 0.663 & MML-EM \\
PSN-IRT (ref.) & --- & --- & Neural \\
\bottomrule
\end{tabular}
\caption{IRT refit parameter recovery. MML-EM fails to recover discrimination ($a_r \ll 0.2$).}
\label{tab:irt_refit}
\end{table}

%-------------------------------------------------------------------
\section{Extended Analysis}
\label{sec:extended}
%-------------------------------------------------------------------

\subsection{Power Function (Seven Benchmarks)}

Adding BBH (6,511 items, power = 0.86) and MATH (5,000 items, power = 0.74) improves the power function fit to $R^2 = 0.82$ (from 0.65 with five benchmarks).
The minimum for 80\% power at 15\% contamination is approximately 5,500 items.

\begin{table}[h]
\centering
\small
\begin{tabular}{lrcc}
\toprule
Benchmark & $n_{\text{items}}$ & Power (sim.) & Power (2-sided) \\
\midrule
HumanEval & 164 & --- & --- \\
ARC-C & 1,172 & --- & --- \\
GSM8K & 1,319 & --- & --- \\
MATH & 5,000 & 0.74 & 0.82 \\
BBH & 6,511 & 0.86 & 0.91 \\
HellaSwag & 10,042 & --- & --- \\
MMLU & 14,042 & 0.979 & --- \\
\bottomrule
\end{tabular}
\caption{Power function across seven benchmarks ($R^2 = 0.82$, 15\% contamination).}
\label{tab:power_function}
\end{table}

\subsection{Multi-Benchmark Detection}

GSM8K and HellaSwag show significant but reduced effects: GSM8K (Qwen $|d| = 1.85$, Llama $|d| = 1.62$), HellaSwag (Qwen $|d| = 2.10$, Llama $|d| = 1.90$), consistent with the power function.

\begin{table}[h]
\centering
\small
\begin{tabular}{lccc}
\toprule
Model $\times$ Benchmark & $|d|$ & MF \\
\midrule
Qwen $\times$ MMLU & 3.15 & .38 \\
Qwen $\times$ GSM8K & 1.85 & --- \\
Qwen $\times$ HellaSwag & 2.10 & --- \\
Llama $\times$ MMLU & 2.72 & .35 \\
Llama $\times$ GSM8K & 1.62 & --- \\
Llama $\times$ HellaSwag & 1.90 & --- \\
\bottomrule
\end{tabular}
\caption{Multi-benchmark results at 15\% contamination.}
\label{tab:multi_bench}
\end{table}

\subsection{Natural Experiment: Base-to-Chat SFT}

Among 9 complete base--chat pairs from Open LLM Leaderboard v1, HellaSwag shows significant negative $\Delta \ell z^*$ (mean $= -3.4$, $t = -3.97$, $p = 0.001$) and GSM8K shows significant positive shift (mean $= +3.92$, $t = 3.35$, $p = 0.009$).
MMLU ($p = 0.15$) and ARC-C ($p = 0.38$) are non-significant.
This demonstrates $\ell z^*$ sensitivity to SFT-induced changes without ground truth.

\begin{table}[h]
\centering
\small
\begin{tabular}{lcccc}
\toprule
Benchmark & Mean $\Delta \ell z^*$ & $t$ & $p$ \\
\midrule
HellaSwag & $-$3.40 & $-$3.97 & 0.001 \\
GSM8K & +3.92 & 3.35 & 0.009 \\
MMLU & --- & --- & 0.15 \\
ARC-C & --- & --- & 0.38 \\
\bottomrule
\end{tabular}
\caption{Base-to-chat $\Delta \ell z^*$ ($n = 9$ pairs). HellaSwag and GSM8K show significant SFT effects.}
\label{tab:natural_exp}
\end{table}

%-------------------------------------------------------------------
\section{Related Work}
\label{sec:related}
%-------------------------------------------------------------------

\paragraph{Contamination detection.}
Methods span five families: (1) time-based partitioning; (2) MIA/Min-K\%++ operating on log-probabilities; (3) probing methods (Time Travel) requiring interactive access; (4) performance-based testing (ConStat); (5) in-context learning effects (CoDeC).
All require more than binary responses.
Wang et al.\ (2025) show all MIA methods fail post-RL.

\paragraph{IRT for LLMs.}
Several works apply IRT for ability estimation and benchmark analysis, but none use person-fit for contamination detection.
PSN-IRT provides our validated 4PL parameters.

\paragraph{Person-fit statistics.}
The psychometric person-fit literature spans five decades, with $\ell z^*$ as the most widely used statistic.
Our work is the first application to LLM benchmark contamination.

%-------------------------------------------------------------------
\section{Discussion and Conclusion}
\label{sec:conclusion}
%-------------------------------------------------------------------

IRT-FIT establishes person-fit statistics as a principled contamination detection method at the minimal-information operating point.
Using only binary response patterns, it achieves $|d| > 2.0$ across four model families, $8\times$ the effect size of classical indices, and uniquely retains detection power post-RL ($|d| = 2.1$).
The $\theta$-fixed methodology universally restores monotonic dose-response, and the memorization fraction (33--38\%) bridges theoretical and observed power.

\paragraph{Limitations.}
IRT-FIT requires pre-calibrated 4PL parameters (currently available via PSN-IRT).
MF depends on the contamination mechanism (QLoRA SFT here) and may vary with other methods.
Item-level detection remains near-random for all binary-only methods.
Small benchmarks ($n < 5{,}500$) have insufficient power at low contamination.
Our scope targets intentional SFT contamination; pretraining exposure produces smooth distributional shifts outside this framework.

\paragraph{Future work.}
13,000+ per-item v1 records are available for large-scale audits with improved neural IRT estimation.
Cross-mechanism validation and generative benchmark extension are natural next steps.

\paragraph{Domain-level analysis.}
Table~\ref{tab:domain_lz} in the appendix presents domain-level $\ell z^*$ for selected models, revealing contamination typology differences (uniform for Gemini-1.5, selective for Doubao-Pro).

\appendix
\section{Supplementary Tables}

\begin{table}[h]
\centering
\small
\begin{tabular}{lccccc}
\toprule
Model & STEM & Humanities & Soc.\ Sci. & Other & Overall \\
\midrule
Gemini-1.5 & 21.2 & 17.9 & 15.9 & 15.3 & 31.81 \\
Doubao-Pro & $-$2 & $-$13 & --- & --- & --- \\
Hunyuan & \multicolumn{4}{c}{\textit{uneven}} & --- \\
\bottomrule
\end{tabular}
\caption{Domain-level $\ell z^*$ (MMLU). Gemini-1.5: uniform; Doubao-Pro: selective.}
\label{tab:domain_lz}
\end{table}

\begin{table}[h]
\centering
\small
\begin{tabular}{lccccc}
\toprule
Model & MMLU $\ell z^*$ & HSwag & GSM8K & ARC-C & $\ell z_m$ \\
\midrule
Gemini-1.5 & 31.81 & --- & --- & --- & 48.5 \\
Clean mean & --- & --- & --- & --- & 28.4$\pm$12.6 \\
Ctrl-contam & --- & --- & --- & --- & 12--15 \\
\bottomrule
\end{tabular}
\caption{Multiscale $\ell z_m$ rankings. Cross-benchmark aggregation provides 15\% sensitivity gain.}
\label{tab:lzm}
\end{table}

\end{document}
